%==============================================================================
% Research Methods - Group assignment - proposals
%==============================================================================

% This loads the article template. For a text written in English, add the
% "english" option, i.e. \documentclass[english]{hogent-article}
\documentclass{hogent-article}

% Include bibliography file
\addbibresource{references.bib}

% Info about the assignment, course and study programme

\studyprogramme{Professionele bachelor toegepaste informatica}
\course{Research Methods, 2E1}
\assignmenttype{Groepsopdracht}
\academicyear{2025-2026}
\title{AI-bias bij geautomatiseerde aanwerving: naar verantwoord gebruik van AI-screeningplatformen}

\author{Simon Borghart}
\email{simon.borghart@student.hogent.be}

\author{Emiel De Pelsmaeker}
\email{emiel.depelsmaeker@student.hogent.be}

\author{Simon Boey}
\email{simon.boey@student.hogent.be}

\projectrepo{https://github.com/SaZoMi/RM-Group-64-interprofessionele.git}

\specialisation{AI \& Data Engineering}

\keywords{algoritmische aanwerving, AI-bias, discriminatie, bias-audit, geautomatiseerde screening}

\begin{document}

\begin{abstract}
  Werkgevers zetten in toenemende mate AI-gedreven platformen in om cv's en
  sollicitaties automatisch te screenen en te rangschikken. In de Verenigde
  Staten lopen momenteel twee collectieve rechtszaken tegen aanbieders van
  dergelijke platformen: Mobley v. Workday, waarin een federale rechtbank
  oordeelde dat het aanwervingsplatform van Workday kan optreden als agent
  van de werkgever en zo onderworpen is aan de federale
  leeftijdsdiscriminatiewet, en Kistler et al. v. Eightfold AI Inc., waarin
  gesteld wordt dat het platform kandidaten scoorde en uitfilterde zonder de
  wettelijk verplichte kennisgevingen. Beide zaken illustreren een breder
  probleem dat in peer-reviewed onderzoek is aangetoond: algoritmische
  aanwervingssystemen kunnen bestaande vooroordelen op basis van leeftijd,
  ras en geslacht reproduceren of versterken. Dit onderzoeksvoorstel
  onderzoekt welke combinatie van technische en organisatorische maatregelen
  het meest geschikt is voor een werkgever die een AI-aanwervingsplatform
  zoals dat van Workday of Eightfold AI inzet, om discriminatie van
  sollicitanten te beperken zonder de efficiëntiewinst van geautomatiseerde
  screening significant te verminderen.
\end{abstract}

\tableofcontents

\section{Inleiding}%
\label{sec:inleiding}

In mei 2025 liet een federale rechtbank in Californië een collectieve
rechtszaak toe tegen softwarebedrijf Workday. Eiser Derek Mobley
solliciteerde over een periode van zeven jaar meer dan honderd keer via het
aanwervingsplatform van Workday en werd telkens binnen enkele minuten
afgewezen. De rechtbank oordeelde dat het platform kan optreden als "agent"
van de werkgevers die het gebruiken, waardoor Workday zelf aansprakelijk kan
worden gesteld onder de federale leeftijdsdiscriminatiewet (ADEA). De
vordering kan mogelijk miljoenen sollicitanten ouder dan veertig jaar
omvatten~\autocite{JonesWalker2026}. In januari 2026 volgde een vergelijkbare
zaak tegen Eightfold AI, een aanwervingsplatform dat onder meer door
Microsoft, PayPal, Morgan Stanley, Starbucks, Chevron en Bayer wordt
gebruikt. De vordering, aangespannen onder de federale Fair Credit Reporting
Act (FCRA), stelt dat Eightfold AI data van meer dan een miljard werknemers
verzamelde, kandidaten scoorde op een schaal van nul tot vijf, en laag
scorende kandidaten uitfilterde vóór menselijke beoordeling, zonder de
wettelijk verplichte kennisgevingen~\autocite{JonesWalker2026}. Jones
Walker LLP is een advocatenkantoor gespecialiseerd in arbeidsrecht en biedt
in die hoedanigheid een praktijkgerichte, feitelijk betrouwbare beschrijving
van het verloop van beide rechtszaken. De bredere argumentatie in dit
onderzoeksvoorstel over de oorzaken en de aanpak van AI-bias steunt echter
op de peer-reviewed bronnen~\autocite{Fabris2025,Chen2023}, niet op de
juridisch-strategische duiding van het kantoor zelf. Beide zaken worden ook
beschreven in een analyse van de University of Miami Law Review, die erop
wijst dat de uitspraak in de zaak tegen Workday de juridische grenzen van
AI-gedreven aanwerving mee zal bepalen, en dat de zaak tegen Eightfold AI
een bijkomende juridische vraag stelt, namelijk of AI-aanwervingstools
onder de Fair Credit Reporting Act vallen~\autocite{Adams2026}.

Deze twee zaken zijn geen alleenstaande incidenten. Fabris et al. (2025)
brengen in een peer-reviewed, multidisciplinair overzicht in kaart hoe
technische, juridische en organisatorische factoren samen bijdragen aan bias
in algoritmische aanwervingssystemen~\autocite{Fabris2025}. Chen (2023)
bevestigt in een peer-reviewed literatuuronderzoek en bevraging dat
AI-gedreven recrutering leidt tot discriminatoire aanwervingspraktijken op
basis van onder meer geslacht, ras en persoonlijkheidskenmerken, en wijst
erop dat deze bias vaak voortkomt uit beperkte of onevenwichtige
trainingsdata en uit vooroordelen bij de ontwerpers van de
algoritmen~\autocite{Chen2023}. Recent empirisch onderzoek, nog niet in een
tijdschrift gepubliceerd maar wel publiek beschikbaar en methodologisch
navolgbaar, toont bovendien dat de bias zich op meerdere manieren
manifesteert: Wen et al. (2025) meten met een eigen benchmark aanzienlijke
raciale en gender-bias in de manier waarop grote taalmodellen cv's
beoordelen~\autocite{Wen2025}, terwijl Xu et al. (2025) aantonen dat
taalmodellen die gebruikt worden om cv's te screenen, systematisch de
voorkeur geven aan cv's die door datzelfde model werden opgesteld, met een
zelfvoorkeur van 67 tot 82 procent en een 23 tot 60 procent hogere kans op
shortlisting~\autocite{Xu2025}.

Dit onderzoeksvoorstel richt zich op HR- en compliance-verantwoordelijken bij
middelgrote tot grote organisaties die een AI-aanwervingsplatform zoals dat
van Workday of Eightfold AI inzetten of overwegen in te zetten, en die
verantwoordelijk zijn voor het beperken van de juridische en
reputatierisico's rond geautomatiseerde aanwerving. Dit is dus niet
"HR-professionals" in het algemeen, maar specifiek de teams die de
leverancierskeuze en het beleid rond geautomatiseerde screening bepalen.

Op basis van dit concrete, benoemde probleem, verankerd in de rechtszaken
tegen Workday en Eightfold AI, luidt de hoofdonderzoeksvraag van dit
onderzoeksvoorstel:

\begin{quote}
  Welke combinatie van technische en organisatorische maatregelen
  (bias-audits, transparante scoringscriteria en menselijk toezicht op
  automatische afwijzingen) is het meest geschikt voor een werkgever die een
  AI-aanwervingsplatform zoals dat van Workday of Eightfold AI inzet, om
  discriminatie van sollicitanten te beperken zonder de efficiëntiewinst van
  geautomatiseerde screening significant te verminderen?
\end{quote}

Deze hoofdvraag is een open vraag, geen ja-neevraag en geen opzoekvraag, en
wordt beantwoord vanuit de concrete situatie van een werkgever die een
platform zoals dat van Workday of Eightfold AI inzet. Om deze hoofdvraag te
beantwoorden, wordt de vraag opgesplitst in deelvragen voor zowel het
probleemdomein als het oplossingsdomein.

Deelvragen probleemdomein:
\begin{itemize}
  \item Welke technische oorzaken liggen aan de basis van bias in
    AI-gedreven cv-screening en rangschikkingssystemen zoals die van
    Workday en Eightfold AI?
  \item Welke juridische risico's lopen werkgevers die dergelijke systemen
    inzetten, in het licht van regelgeving zoals de ADEA en de FCRA?
  \item In welke mate zijn sollicitanten zich bewust van de rol van AI in
    het aanwervingsproces en welke impact heeft dit op hun vertrouwen in de
    werkgever?
\end{itemize}

Deelvragen oplossingsdomein:
\begin{itemize}
  \item Hoe werkt een bias-audit van een AI-aanwervingssysteem in de
    praktijk en welke meetmethoden worden hiervoor gebruikt?
  \item Welke vormen van menselijk toezicht op automatische
    afwijzingsbeslissingen worden in de literatuur voorgesteld en hoe
    verhouden deze zich tot de efficiëntiewinst van automatisering?
  \item Hoe kan een werkgever transparantie over scoringscriteria
    organiseren zonder de onderliggende modellen volledig openbaar te
    maken?
  \item Hoe kan de impact van deze maatregelen op zowel bias-reductie als
    efficiëntie gemeten en vergeleken worden?
\end{itemize}

Het doel van dit onderzoek is een afwegingskader op te stellen dat een
werkgever, primair gericht op een organisatie die een platform zoals dat van
Workday of Eightfold AI inzet, en bij uitbreiding op vergelijkbare
werkgevers, helpt om AI-gedreven aanwervingssystemen op een verantwoorde
manier in te zetten. Het onderzoek is specifiek doordat het zich beperkt tot
drie met naam benoemde maatregelen binnen één concreet probleem, meetbaar
doordat bias-reductie, doorlooptijd en het aantal terechte of onterechte
afwijzingen kwantificeerbaar zijn, haalbaar doordat het steunt op publiek
beschikbare datasets en gepubliceerde rechtszaakdocumenten zonder interne
bedrijfsdata te vereisen, relevant doordat het probleem recent is aangetoond
in peer-reviewed onderzoek en zich voordoet bij publiek gedocumenteerde
bedrijven, en tijdsgebonden doordat de volledige aanpak binnen twaalf weken
wordt uitgewerkt.

De rest van dit document is als volgt opgebouwd. Sectie~\ref{sec:literatuurstudie}
bespreekt de bestaande literatuur over bias in algoritmische aanwerving.
Sectie~\ref{sec:methodologie} beschrijft de voorgestelde onderzoeksaanpak,
inclusief het technische luik. Sectie~\ref{sec:resultaten} beschrijft de
verwachte resultaten. Sectie~\ref{sec:conclusie} sluit af met een discussie
en conclusie.

\section{Literatuurstudie}%
\label{sec:literatuurstudie}

Fabris et al. (2025) bieden een multidisciplinair overzicht van bias bij
algoritmische aanwerving, waarin technische bronnen van bias (bijvoorbeeld
onevenwichtige trainingsdata), juridische kaders (waaronder
antidiscriminatiewetgeving in verschillende rechtsgebieden) en
organisatorische factoren (zoals gebrek aan menselijk toezicht) worden
samengebracht~\autocite{Fabris2025}. Dit overzicht vormt het theoretische
kader waarbinnen de rechtszaken tegen Workday en Eightfold AI kunnen worden
begrepen: het toont aan dat de problemen in deze zaken geen alleenstaande
incidenten zijn, maar een structureel patroon volgen dat in de literatuur al
uitgebreid is gedocumenteerd.

Naast deze academische en juridisch-analytische bronnen bestaat er ook een
regelgevend kader dat rechtstreeks van toepassing is. Onder de EU AI Act
worden AI-systemen die gebruikt worden voor de rekrutering of selectie van
kandidaten, waaronder het filteren van sollicitaties en het evalueren van
kandidaten, expliciet geclassificeerd als hoogrisicosystemen (Bijlage III,
punt 4)~\autocite{EUAIAct2024}. Dit betekent dat werkgevers die platformen
zoals Workday of Eightfold AI binnen de Europese Unie inzetten, wettelijk
verplicht zullen zijn tot onder meer risicobeheer, menselijk toezicht en
transparantie, dezelfde soort maatregelen die in dit onderzoeksvoorstel
centraal staan. Deze regelgeving plaatst het probleem dus niet enkel in een
Amerikaanse rechtszaakcontext, maar ook in een Europese regelgevende context
die rechtstreeks relevant is voor de beoogde doelgroep.

Ook in de Verenigde Staten, waar beide rechtszaken tegen Workday en
Eightfold AI gevoerd worden, bestaat officiële overheidsduiding over dit
probleem. De Equal Employment Opportunity Commission (EEOC), de federale
instantie die arbeidsdiscriminatiewetgeving handhaaft, wijst er publiek op
dat bestaande federale antidiscriminatiewetten onverkort van toepassing
blijven wanneer AI wordt ingezet bij werving, screening of evaluatie van
sollicitanten, ongeacht of dit via een eigen systeem of via een externe
leverancier gebeurt~\autocite{EEOC2024}. Dit bevestigt vanuit
overheidsperspectief dat de aansprakelijkheid van de werkgever, zoals die
ook in de zaak Mobley v. Workday naar voor komt, geen geïsoleerde
rechterlijke interpretatie is, maar aansluit bij de bredere handhavingsvisie
van de bevoegde federale instantie. De EEOC publiceerde in mei 2023
uitgebreidere technische richtlijnen specifiek over AI bij aanwerving, maar
deze werden op 27 januari 2025 door de EEOC zelf van de website gehaald na
een federale beleidswijziging. Het hier aangehaalde document uit april 2024
blijft wel publiek beschikbaar en weerspiegelt de blijvend geldende
wettelijke bescherming.

Chen (2023) onderzoekt via literatuuronderzoek en bevraging hoe AI-gedreven
recrutering leidt tot discriminatoire aanwervingspraktijken, en wijst op
twee hoofdoorzaken: beperkte of onevenwichtige trainingsdatasets en
vooroordelen bij de ontwerpers van de algoritmen. Het onderzoek beveelt zowel
technische maatregelen aan, zoals onbevooroordeelde datasetkaders en
verbeterde algoritmische transparantie, als beheersmaatregelen, zoals interne
ethische governance en extern toezicht~\autocite{Chen2023}. Deze aanbevelingen
vormen een directe basis voor de drie maatregelen die in dit
onderzoeksvoorstel centraal staan: bias-audits, transparante
scoringscriteria en menselijk toezicht.

Recent empirisch onderzoek maakt de aard van deze bias concreter meetbaar.
Wen et al. (2025) ontwikkelen de FAIRE-benchmark om raciale en gender-bias
te meten in de manier waarop grote taalmodellen cv's beoordelen over
verschillende sectoren heen, met behulp van zowel directe scoring als
rangschikking. Zij stellen vast dat elk onderzocht model een zekere mate van
bias vertoont, al verschillen omvang en richting sterk tussen
modellen~\autocite{Wen2025}. Xu et al. (2025) onderzoeken een ander type
bias: zelfvoorkeur. Via een grootschalig gecontroleerd experiment met
cv-vergelijkingen stellen zij vast dat taalmodellen consistent de voorkeur
geven aan cv's die door datzelfde model werden geschreven of aangepast, ook
wanneer de inhoudelijke kwaliteit gelijk gehouden wordt. Gesimuleerde
aanwervingspijplijnen over vierentwintig beroepen tonen aan dat kandidaten
die hetzelfde taalmodel gebruiken als de beoordelaar, 23 tot 60 procent meer
kans hebben om geselecteerd te worden dan even gekwalificeerde kandidaten met
een door een mens geschreven cv. Zij tonen ook aan dat deze bias met meer
dan de helft kan worden teruggedrongen via eenvoudige interventies die
inspelen op het zelfherkenningsvermogen van het model~\autocite{Xu2025}.

De juridische analyse van Jones Walker LLP (2026) plaatst deze technische
bevindingen in een concrete, actuele bedrijfscontext. De zaak Mobley
v. Workday toont aan dat een aanwervingsplatform juridisch kan worden
beschouwd als agent van de werkgever, met aansprakelijkheid onder de ADEA
tot gevolg. De zaak Kistler et al. v. Eightfold AI Inc. toont aan dat het
ontbreken van wettelijk verplichte kennisgevingen bij geautomatiseerde
scoring tot aansprakelijkheid onder de FCRA kan leiden, en dat dit risico
zich voordoet bij platformen die door grote, bekende werkgevers worden
ingezet~\autocite{JonesWalker2026}. Samen bevestigen deze bronnen zowel het
technische probleem (bias in de onderliggende modellen), het juridische
risico (aansprakelijkheid onder bestaande regelgeving), als de
organisatorische oplossingsrichting (audits, transparantie, toezicht) die in
dit onderzoeksvoorstel worden samengebracht.

\section{Methodologie}%
\label{sec:methodologie}

Het onderzoek wordt uitgevoerd in vier fasen, elk twee tot vier weken,
samen passend binnen de beschikbare twaalf weken van dit onderzoekstraject.

\textbf{Fase 1, probleemdomein (weken 1 tot 3).} Via literatuuronderzoek en
analyse van de publiek beschikbare rechtszaakdocumenten van Mobley
v. Workday en Kistler et al. v. Eightfold AI Inc. wordt uitgezocht welke
technische en juridische oorzaken aan de basis liggen van bias bij
AI-gedreven cv-screening. Deze fase levert een probleembeschrijving op die
zowel de technische oorzaken van bias als het juridische risicokader
(ADEA, FCRA) in kaart brengt.

\textbf{Fase 2, oplossingsdomein vanuit de literatuur (weken 4 tot 6).} De
drie voorgestelde maatregelen, bias-audits, transparante scoringscriteria
en menselijk toezicht op automatische afwijzingen, worden beschreven en
onderling vergeleken op basis van de bestaande literatuur, met bijzondere
aandacht voor de methoden die Wen et al. (2025) en Xu et al. (2025)
gebruiken om bias te meten. Deze fase levert een literatuurgebaseerd
afwegingskader op.

\textbf{Fase 3, technisch en praktisch luik (weken 7 tot 10).} Er wordt een
eenvoudig bias-auditscript geïmplementeerd en toegepast op een publiek
beschikbare cv- of aanwervingsdataset. Het script meet, naar analogie met de
methode van Wen et al. (2025), het verschil in scoring of shortlisting-kans
tussen kandidaten met vergelijkbare kwalificaties maar verschillende
gesuggereerde demografische kenmerken. Deze fase levert gemeten
PoC-resultaten op die aantonen in welke mate bias meetbaar en detecteerbaar
is met eenvoudige, voor een werkgever haalbare middelen. Dit vult het
cesuur-criterium in dat een technische component vereist naast een louter
theoretisch of academisch werk.

\textbf{Fase 4, synthese (weken 11 en 12).} De bevindingen uit de
literatuur en de PoC worden samengebracht tot een afwegingskader voor
werkgevers die een AI-aanwervingsplatform zoals dat van Workday of
Eightfold AI inzetten. Deze fase levert het volledige eindrapport op.

Elke fase levert een concreet milestone op: een probleembeschrijving met
juridisch risicokader na fase 1, een literatuurgebaseerd afwegingskader na
fase 2, gemeten PoC-resultaten na fase 3, en het volledige eindrapport na
fase 4. Deze opeenvolging van vier fasen van telkens twee tot vier weken
past binnen de beschikbare twaalf weken van dit onderzoekstraject.

\section{Verwachte resultaten}%
\label{sec:resultaten}

Dit onderzoek verwacht een afwegingskader op te leveren waarmee een
werkgever, primair gericht op een organisatie die een platform zoals dat van
Workday of Eightfold AI inzet, en bij uitbreiding vergelijkbare werkgevers,
kan inschatten welke combinatie van bias-audits, transparante
scoringscriteria en menselijk toezicht het meest geschikt is om
discriminatie te beperken zonder de efficiëntiewinst van geautomatiseerde
screening significant te verminderen. De PoC-resultaten uit fase 3 zullen
concreet aantonen in welke mate bias in een AI-gedreven screeningsysteem
detecteerbaar is met een eenvoudig, voor een werkgever haalbaar auditscript,
en welke orde van grootte van bias daarbij gemeten wordt in vergelijking met
de bevindingen van Wen et al. (2025) en Xu et al. (2025).

\section{Discussie en conclusie}%
\label{sec:conclusie}

De rechtszaken Mobley v. Workday en Kistler et al. v. Eightfold AI Inc.
tonen aan dat bias in AI-gedreven aanwervingssystemen geen theoretisch
probleem is, maar reële juridische en maatschappelijke gevolgen heeft voor
zowel sollicitanten als werkgevers. Dit onderzoeksvoorstel biedt een aanpak
om, vertrekkend vanuit dit concrete probleem, een afwegingskader te
ontwikkelen voor een werkgever zoals een organisatie die een platform zoals
dat van Workday of Eightfold AI inzet. Door literatuuronderzoek te
combineren met een technisch-praktisch luik, levert het onderzoek zowel
theoretisch onderbouwde inzichten als een concrete, meetbare bijdrage over
de detecteerbaarheid van bias in AI-gedreven aanwervingssystemen.

\printbibliography[heading=bibintoc]

\end{document}
